\vspace{-3mm}
\section{Conclusion and limitation}
\vspace{-3mm}
We propose \texttt{Where2comm}, a novel communication-efficient collaborative perception framework. The core idea is to exploit a spatial confidence map at each agent to promote pragmatic compression, assisting agents to decide what to communicate with whom, and whose information to aggregate. Each agent offers spatially sparse, yet perceptually critical features to support other agents; meanwhile, requests complementary information from others in multi-round communication. Comprehensive experiments covering multi-type agents and multi-modality inputs show that \texttt{Where2comm} achieves far superior trade-off between perception performance and communication bandwidth.

\textbf{Limitation and future work.} 
The current work focuses on perceptually critical spatial areas. In future, we plan to expand a similar idea to the temporal dimension and determine critical time stamps. More cost will be reduced by exploring when to communicate. We also expect that more methods on pragmatic compression and emergent communication could be applied to collaborative perception.

\textbf{Acknowledgment.}
This research is partially supported by the National Key R\&D Program of China under Grant 2021ZD0112801, National Natural Science Foundation of China under Grant 62171276, the Science and Technology Commission of Shanghai Municipal under Grant 21511100900, CCF-DiDi GAIA Research Collaboration Plan 202112 and CALT Grant 2021-01.